\documentclass[11pt,a4paper]{article}
\usepackage{fontspec,amsmath,amssymb,amsthm,microtype}
\setmainfont{lmroman10-regular.otf}[BoldFont=lmroman10-bold.otf,ItalicFont=lmroman10-italic.otf,BoldItalicFont=lmroman10-bolditalic.otf]
\setmonofont{lmmono10-regular.otf}[Scale=MatchLowercase]
\usepackage[margin=25mm]{geometry}
\usepackage[unicode,hidelinks]{hyperref}
\setlength{\emergencystretch}{3em}
\setlength{\parskip}{0.4em}
\setlength{\parindent}{0pt}
\allowdisplaybreaks[2]
\raggedbottom
\title{\texorpdfstring{The actual Brownian KIPLD flow and its stationary mixing theorem}{The actual Brownian KIPLD flow and its stationary mixing theorem}}
\author{LangevinLean mathematical exposition}
\date{}
\begin{document}\maketitle
\noindent Source module: \texttt{KIPLMC/\allowbreak{}Continuous.\allowbreak{}lean}. The accompanying HTML pairs this mathematical exposition with the exact Lean statements and records the scope of the formalization.


This module connects the canonical Brownian kinetic diffusion to the physical particle system and proves Proposition 3 for its actual output law. The initial position is an arbitrary probability law with finite second moment. Standard Gaussian velocity and the canonical Brownian source are independent of that position; velocity scaling gives the required physical thermal law. The Gibbs law is a proved stationary target, and the Wasserstein distance is the infimum over all couplings.


\section{\texorpdfstring{Construct the source, the physical flow, and its actual endpoint law}{Construct the source, the physical flow, and its actual endpoint law}}


Set \(c=\sqrt N\). The canonical source contains a standard Gaussian velocity and an independent vector Brownian source on the actual coordinate set \(\mathrm{Fin}(d_\theta)\sqcup(\mathrm{Fin}(N)\times\mathrm{Fin}(d_x))\). Use the proved selected kinetic flow with sum potential \(V_N\), smoothness \(NL\), and friction \(c\gamma\). The physical process is \(D_{c^{-1}}Y_{t/c}\). Position is unchanged, so its endpoint pushforward is exactly the canonical position law at time \(T/c\). This exact equality is the bridge used by the convergence proof.



{\raggedright\noindent Lean declarations: \texttt{kipld\allowbreak{}Clock\allowbreak{}Friction}, \texttt{kipld\allowbreak{}Brownian\allowbreak{}Source}, \texttt{kipld\allowbreak{}Brownian\allowbreak{}Flow}, \texttt{kipld\allowbreak{}Brownian\allowbreak{}Position\allowbreak{}Law}, \texttt{kipld\allowbreak{}Brownian\allowbreak{}Position\allowbreak{}Law\_\allowbreak{}eq\_\allowbreak{}canonical}.\par}

\section{\texorpdfstring{Verify the independent thermal initialization}{Verify the independent thermal initialization}}


The complete source law theorem supplies the actual initial phase law \(\mathcal N(0,I)\otimes\nu\). Push it through velocity multiplication by \(N^{-1/2}\). The result is precisely \texttt{velocity\allowbreak{}Target.\allowbreak{}prod \ensuremath{\nu}}; the velocity and position remain independent. The selected flow’s time-zero identity and the almost-everywhere scalar-multiplication rule identify the actual L\ensuremath{{}^{2}} representatives with this pushforward. No initialization-law hypothesis is imposed on the final theorem.



{\raggedright\noindent Lean declarations: \texttt{kipld\allowbreak{}Brownian\allowbreak{}Flow\_\allowbreak{}initial\_\allowbreak{}has\allowbreak{}Law}.\par}

\section{\texorpdfstring{Cancel the clock factors and prove stationary mixing}{Cancel the clock factors and prove stationary mixing}}


The lifted potential has curvature \(Nm\) and smoothness \(NL\). Multiplication by \(N\) transforms \(m+L\le\gamma^2\) into \(Nm+NL\le(\sqrt N\gamma)^2\), so the actual canonical Gibbs mixing theorem applies. Its exponent is \(-(Nm)/(\sqrt N\gamma)\,(T/\sqrt N)=-mT/\gamma\). The resulting estimate is \(W_2(\nu_T,\pi_Z)\le\sqrt2e^{-mT/\gamma}W_2(\nu,\pi_Z)\). The canonical theorem proves Gibbs stationarity internally. Separate moment and probability theorems certify that all real Wasserstein quantities are finite.



{\raggedright\noindent Lean declarations: \texttt{kipld\_\allowbreak{}scaled\_\allowbreak{}exponent}, \texttt{kipld\allowbreak{}Brownian\allowbreak{}Position\allowbreak{}Law\_\allowbreak{}wasserstein\_\allowbreak{}bound}, \texttt{kipld\allowbreak{}Brownian\allowbreak{}Position\allowbreak{}Law\_\allowbreak{}mem\allowbreak{}Lp}, \texttt{kipld\allowbreak{}Brownian\allowbreak{}Position\allowbreak{}Law\_\allowbreak{}probability}, \texttt{kipld\allowbreak{}Brownian\allowbreak{}Position\allowbreak{}Law\_\allowbreak{}wasserstein2\_\allowbreak{}ne\_\allowbreak{}top}.\par}

\section{\texorpdfstring{Identify the actual physical force in the integral equation}{Identify the actual physical force in the integral equation}}


The drift lifted to L\ensuremath{{}^{2}} represents pointwise application of the true sum gradient. The almost-everywhere scalar multiplication and composition identities show that \(N^{-1}\operatorname{lpDrift}(\nabla V_N)=\operatorname{lpDrift}(F)\), where \(F=\nabla V_N/N\) is the previously verified particle force. Apply the exact kinetic clock theorem to the selected canonical solution. The result is the full physical-time integral equation with friction \(\gamma\), the actual force \(F\), and input \(N^{-1/2}B_{t/\sqrt N}\). The shared Brownian scaling theorem identifies its thermal amplitude as \(\sqrt{2\gamma/N}\).



{\raggedright\noindent Lean declarations: \texttt{kipld\_\allowbreak{}lp\allowbreak{}Force\_\allowbreak{}clock}, \texttt{kipld\allowbreak{}Brownian\allowbreak{}Flow\_\allowbreak{}integral}.\par}

\section{\texorpdfstring{Project to the parameter and add the statistical error}{Project to the parameter and add the statistical error}}


The parameter projection has operator norm at most one, so pushing any coupling through it cannot increase its quadratic cost. Taking the true coupling infimum bounds the parameter Wasserstein distance by the full position distance. The initial independent coupling \(\nu\otimes\pi_Z\) bounds that infimum by \((\mathbb E\|Z_0-\bar Z_*\|^2)^{1/2}\), giving precisely the printed proposition. A second endpoint adds the already proved sharp equilibrium concentration radius \(\sqrt{d_\theta/(Nm)}\) by the L\ensuremath{{}^{2}} triangle inequality, giving a direct parameter RMS bound about the unique marginal-likelihood maximizer.



{\raggedright\noindent Lean declarations: \texttt{kipld\_\allowbreak{}parameter\_\allowbreak{}sampling\_\allowbreak{}bound}, \texttt{kipld\_\allowbreak{}parameter\_\allowbreak{}RMS\_\allowbreak{}bound}.\par}

\end{document}
